/home/sybil/.tex_templates/template.tex
\setupHW{MATH 341: Howework 5}
\begin{document}
\maketitle
\section*{Exercise 1 (Nonlinear least squares)}
Consider the data and the power-law model
\[
	t = \begin{bmatrix} 1 \\ 2 \\ 3 \\ 4 \\ 5 \end{bmatrix}, \quad
	y = \begin{bmatrix} 1.2 \\ 2.8 \\ 5.1 \\ 8.0 \\ 14.7 \end{bmatrix},
\]
\[ f(t; c) = c_1 t^{c_2}, \quad c = \begin{bmatrix} c_1 \\ c_2 \end{bmatrix}. \]

\subsection*{Part (a)}
Define the nonlinear least-squares problem for fitting the model to the data. In particular, write down the residuals and the residual vector $r(c)$.
\[ r_i(c) = c_1 t_i^{c_2} - y_i, \quad i = 1, \dots, 5 \]

\[
	\begin{gathered}
		\text{Consider the form below: } \\
		\begin{bmatrix}
			t_1    & \cdots & \cdots & 0      \\
			0      & t_2    &        & \vdots \\
			\vdots &        & \ddots & 0      \\
			0      & \cdots & \cdots & t_5    \\
		\end{bmatrix}^{c_2}
		\begin{bmatrix}
			c_1    \\
			c_1    \\
			\vdots \\
			c_1    \\
		\end{bmatrix} =
		\begin{bmatrix}
			\hat{y}_1 \\
			\hat{y}_2 \\
			\vdots    \\
			\hat{y}_N \\
		\end{bmatrix}
	\end{gathered}
\]

\[
	\implies \vec{r} = \vec{\hat{y}}-\vec{y}
\]


\subsection*{Part (b)}
Define: $E(c) = \frac{1}{2} \|r(c)\|_2^2.$
Set $F(c) = \nabla E(c)$, where $\nabla E(c)$ is regarded as a row vector. Reformulate the stationarity condition as where $F(c)^T$ is a column vector.
\[ F(c)^T = 0, \]
Show that where $J(c) = D r(c)$, and compute the Jacobian matrix $J(c)$ explicitly.
\[ F(c)^T = J(c)^T r(c), \]

% Your solution here
\[
	E(c) = \frac{1}{2} \|r(c)\|_2^2 ~\equiv~ \frac{1}{2} \vec{r}^{\,T} \vec{r} \implies \nabla E :=
	\left(\frac{\partial \vec{r}}{\partial c}\right)^{T} \vec{r} = J(c)^{T} \vec{r} = F^T(c)
\]

\[
	\therefore ~F(c)^T = 0 \iff J(c)r(c)=0
\]

\[
	\begin{gathered}
		c_1t_i^{c_2} \frac{\partial r_i}{\partial c_1} = t_i^{c_2} \quad
		c_1t_i^{c_2} \frac{\partial r_i}{\partial c_2} = c_1t_i^{c_2}log(t_i)
	\end{gathered}
	\implies
	J :=
	\begin{bmatrix}
		t_1^{c_2} & c_1t_1^{c_2}log(t_1) \\
		\vdots    & \vdots               \\
		t_5^{c_2} & c_1t_5^{c_2}log(t_5) \\
	\end{bmatrix}
\]

\subsection*{Part (c/d)}
Formulate the Gauss-Newton algorithm for this problem, starting from the initial guess
\[ c^{(0)} = \begin{bmatrix} 1 \\ 1 \end{bmatrix}. \]
As a hint, if $A = J(c^{(k)})$, then the Gauss-Newton step $\delta c$ is defined by
\[ A^T A \delta c = -A^T r(c^{(k)}) \]
\[ c^{(k+1)} = c^{(k)} + \delta c. \]

If possible, implement the algorithm in your favorite language (MATLAB, Python, Julia, C/C++, \dots). Report the approximate fitted parameters and, if you wish, include a plot of the data together with the fitted curve.

\begin{cbox}{catred}{Code/Algorithim}
	cbox\inputminted[linenos, breaklines, frame=lines, framesep=2mm, numbersep=8pt, tabsize=4]{python}{/home/sybil/Documents/School/2026-Spring-bk/MATH-341/HW/06/code/main.py}
\end{cbox}

\includegraphics[width=\textwidth]{/home/sybil/Documents/School/2026-Spring-bk/MATH-341/HW/06/code/gauss_newton_exponential_fit.png}

\begin{table}[h]
	\centering
	\begin{tabular}{l l}
		\hline
		Parameter set            & Values                     \\
		\hline
		Initial guess $c^{(0)}$  & $[1.0,\ 1.0]$              \\
		Optimized parameters $c$ & $[0.96590464,\ 0.5465158]$ \\
		\hline
	\end{tabular}
\end{table}

\section*{Exercise 2 (Linearization by taking logarithms)}

Use the same data and model as in Problem 1, but now fit the model through logarithmic linearization.

\subsection*{Part (a)}
Starting from
\[ y_i = c_1 t_i^{c_2}, \qquad i=1,\dots,N, \]
and using the assumptions $y_i > 0$ and $t_i > 0$, take logarithms of both sides:
\[ \log y_i = \log c_1 + c_2 \log t_i. \]
With the substitutions
\[ z_i = \log t_i, \qquad w_i = \log y_i, \qquad \alpha = \log c_1, \qquad \beta = c_2, \]
the model becomes
\[ w_i = \alpha + \beta z_i, \qquad i=1,\dots,N. \]
Therefore, the fitting problem is to find $\alpha$ and $\beta$ that minimize
\[ \sum_{i=1}^N \bigl(w_i - (\alpha + \beta z_i)\bigr)^2. \]
This is a linear least-squares problem in the unknown vector
\[ u = \begin{bmatrix} \alpha \\ \beta \end{bmatrix}. \]
In matrix form,
\[ Xu \approx w, \]
where
\[
	X =
	\begin{bmatrix}
		1      & z_1    \\
		1      & z_2    \\
		\vdots & \vdots \\
		1      & z_N
	\end{bmatrix},
	\qquad
	u =
	\begin{bmatrix}
		\alpha \\
		\beta
	\end{bmatrix},
	\qquad
	w =
	\begin{bmatrix}
		w_1    \\
		w_2    \\
		\vdots \\
		w_N
	\end{bmatrix}.
\]
Thus, the linear least-squares problem is
\[
	\min_{u \in \mathbb{R}^2} \|Xu - w\|_2^2.
\]
Comparing with the generic form, we have
\[
	x_{i,1} = 1, \qquad x_{i,2} = z_i, \qquad \theta_1 = \alpha, \qquad \theta_2 = \beta, \qquad \hat{y}_i = w_i.
\]
After solving for $u$, recover the original parameters by
\[
	c_1 = e^{\alpha}, \qquad c_2 = \beta.
\]


\subsection*{Part (b)}
Define
\[
	E_{\log}(u)=\frac{1}{2}\|Xu-w\|_2^2=\frac{1}{2}(Xu-w)^T(Xu-w).
\]
Then
\[
	F_{\log}(u)=\nabla E_{\log}(u)=X^T(Xu-w)=X^TXu-X^Tw.
\]
Hence the stationarity condition is
\[
	F_{\log}(u)^T=0
	\iff F_{\log}(u)=0
	\iff X^T(Xu-w)=0
	\iff X^TXu=X^Tw.
\]

Assuming the observations are independent, the likelihood is
\[
	J(\theta)=\prod_{i=1}^{N}\big[\alpha(\theta^Tx_i)\big]^{y_i}\big[1-\alpha(\theta^Tx_i)\big]^{1-y_i}.
\]
Taking logs gives the log-likelihood
\[
	\ell(\theta)=\log J(\theta)=\sum_{i=1}^{N}\Big[y_i\log\alpha(\theta^Tx_i)+(1-y_i)\log\big(1-\alpha(\theta^Tx_i)\big)\Big].
\]

For a single sample $(x_i,y_i)$, let
\[
	\ell^{(i)}(\theta)=y_i\log\hat{y}_i+(1-y_i)\log(1-\hat{y}_i),
	\qquad
	\hat{y}_i=\alpha(\theta^Tx_i).
\]
Using the chain rule,
\[
	\frac{\partial \ell^{(i)}(\theta)}{\partial\theta_j}
	=
	\frac{\partial \ell^{(i)}}{\partial \hat{y}_i}
	\cdot
	\frac{\partial \hat{y}_i}{\partial (\theta^Tx_i)}
	\cdot
	\frac{\partial (\theta^Tx_i)}{\partial \theta_j}.
\]
Now,
\[
	\frac{\partial \ell^{(i)}}{\partial \hat{y}_i}
	=
	\frac{y_i}{\hat{y}_i}-\frac{1-y_i}{1-\hat{y}_i},
	\qquad
	\frac{\partial \hat{y}_i}{\partial (\theta^Tx_i)}=\hat{y}_i(1-\hat{y}_i),
	\qquad
	\frac{\partial (\theta^Tx_i)}{\partial \theta_j}=x_{ij}.
\]
Therefore,
\[
	\frac{\partial \ell^{(i)}(\theta)}{\partial\theta_j}
	=
	\left(\frac{y_i}{\hat{y}_i}-\frac{1-y_i}{1-\hat{y}_i}\right)\hat{y}_i(1-\hat{y}_i)x_{ij}
	=
	(y_i-\hat{y}_i)x_{ij}.
\]
Equivalently, in vector form,
\[
	\nabla \ell^{(i)}(\theta)=(y_i-\hat{y}_i)x_i.
\]

If instead $J^{(i)}(\theta)$ denotes the \emph{negative} log-likelihood (the cost minimized in logistic regression), then
\[
	\frac{\partial J^{(i)}(\theta)}{\partial\theta_j}=(\hat{y}_i-y_i)x_{ij},
	\qquad
	\nabla J^{(i)}(\theta)=(\hat{y}_i-y_i)x_i.
\]

% Your original final line had the right idea, but the sign depends on whether you use
% log-likelihood or negative log-likelihood, and the correct vector form is
% (\hat{y}_i-y_i)x_i or (y_i-\hat{y}_i)x_i.


\subsection*{Part (c)}
Recover the original parameters from
\[ c_1 = e^\alpha, \quad c_2 = \beta. \]
If possible, compare the fitted parameters from Problems 1 and 2 and explain briefly why the two answers need not agree.

% Your solution here

\subsection*{Part (d)}
If possible, implement the log-linearized approach in your favorite language and compare the resulting fit with the Gauss-Newton fit from Problem 1.

% Your solution here

\vspace{0.5cm}
\hrule
\vspace{0.5cm}

\section*{Exercise 3 (Rank-one and low-rank inverse updates)}

\subsection*{Part (a)}
Let $c, d \in \mathbb{R}^n$ and assume that
\[ 1 + d^T c \neq 0. \]
Verify directly that
\[ (I + c d^T)^{-1} = I - \frac{c d^T}{1 + d^T c}. \]
That is, multiply $(I + c d^T)$ by the proposed inverse and simplify.

\[
\begin{aligned}
	(I + c d^T)\left( I - \frac{c d^T}{1 + d^T c} \right)
	&= I + c d^T - \frac{c d^T}{1 + d^T c} - \frac{c d^T c d^T}{1 + d^T c} \\
	&= I + c d^T - \frac{c d^T}{1 + d^T c} - \frac{(d^T c)c d^T}{1 + d^T c} \\
	&= I + c d^T - \frac{(1 + d^T c)c d^T}{1 + d^T c} \\
	&= I
\end{aligned}
\]

\subsection*{Part (b)}
Let $A \in \mathbb{R}^{n \times n}$ be nonsingular, and let $C, D \in \mathbb{R}^{n \times k}$. Assume that
\[ I + D^T A^{-1} C \]
is nonsingular. Verify the Sherman-Morrison-Woodbury formula
\[ (A + C D^T)^{-1} = A^{-1} - A^{-1} C (I + D^T A^{-1} C)^{-1} D^T A^{-1}. \]
Again, verify this by multiplying $(A + C D^T)$ by the inverse and simplifying.

\[
\begin{aligned}
\text{Let } M &= I + D^T A^{-1} C. \\
(A + C D^T)\left(A^{-1} - A^{-1} C M^{-1} D^T A^{-1}\right)
&= A A^{-1} - A A^{-1} C M^{-1} D^T A^{-1} + C D^T A^{-1} - C D^T A^{-1} C M^{-1} D^T A^{-1} \\
&= I - C M^{-1} D^T A^{-1} + C D^T A^{-1} - C D^T A^{-1} C M^{-1} D^T A^{-1} \\
&= I + C D^T A^{-1} - C\left(M^{-1} + D^T A^{-1} C M^{-1}\right)D^T A^{-1} \\
&= I + C D^T A^{-1} - C\left(I + D^T A^{-1} C\right)M^{-1} D^T A^{-1} \\
&= I + C D^T A^{-1} - C D^T A^{-1} \\
&= I.
\end{aligned}
\]

Hence
\[
(A + C D^T)^{-1} = A^{-1} - A^{-1} C (I + D^T A^{-1} C)^{-1} D^T A^{-1}.
\]

\subsection*{Part (c)}
Explain briefly why part (a) is a special case of part (b).

\[
\text{Let: } A = I, \quad C = c, \quad D = d, \quad k = 1 \quad ~st~ \quad CD^T = cd^T, \quad I + D^T A^{-1}C = 1 + d^T c,
\]
\[
\implies (I + cd^T)^{-1} = I - \frac{cd^T}{1 + d^T c}.
\]

\section*{Exercise 4}

Let $A \in \mathbb{R}^{n \times n}$ be a symmetric positive definite matrix and $b \in \mathbb{R}^n$. Consider the quadratic functional
\[ J(x) = \frac{1}{2} x^T A x - x^T b, \quad x \in \mathbb{R}^n. \]
It is known that minimizing $J$ is equivalent to solving the linear system
\[ A x = b. \]
Verify this by computing the partial derivative of $J$ with respect to $x_k$, and showing that
\[ \frac{\partial J}{\partial x_k}(x) = 0 \]
is exactly the $k$-th equation of the system $A x = b$, for each $k = 1, \dots, n$.
\[ J(x) = \frac{1}{2} \sum_{i=1}^n \sum_{j=1}^n a_{ij} x_i x_j - \sum_{i=1}^n b_i x_i, \]
and use that $A$ is symmetric.

% Solution
\[
\begin{aligned}
\frac{\partial J}{\partial x_k}(x)
&= \frac{\partial}{\partial x_k}\left(\frac{1}{2} \sum_{i=1}^n \sum_{j=1}^n a_{ij} x_i x_j - \sum_{i=1}^n b_i x_i\right) \\
&= \frac{1}{2} \sum_{j=1}^n a_{kj} x_j + \frac{1}{2} \sum_{i=1}^n a_{ik} x_i - b_k.
\end{aligned}
\]
Since $A$ is symmetric, $a_{ik} = a_{ki}$. Therefore,
\[
\frac{1}{2} \sum_{i=1}^n a_{ik} x_i = \frac{1}{2} \sum_{i=1}^n a_{ki} x_i.
\]
Relabeling the index $i$ as $j$, we obtain
\[
\frac{\partial J}{\partial x_k}(x)
= \frac{1}{2} \sum_{j=1}^n a_{kj} x_j + \frac{1}{2} \sum_{j=1}^n a_{kj} x_j - b_k
= \sum_{j=1}^n a_{kj} x_j - b_k.
\]
But
\[
\sum_{j=1}^n a_{kj} x_j = (Ax)_k,
\]
so
\[
\frac{\partial J}{\partial x_k}(x) = (Ax)_k - b_k.
\]
Hence,
\[
\frac{\partial J}{\partial x_k}(x) = 0
\quad \Longleftrightarrow \quad
(Ax)_k = b_k
\quad \Longleftrightarrow \quad
\sum_{j=1}^n a_{kj} x_j = b_k,
\]
which is exactly the $k$-th equation of the linear system $Ax=b$. Since this holds for every $k=1,\dots,n$, minimizing $J$ is equivalent to solving $Ax=b$.

\end{document}
